% Part: propositional-logic
% Chapter: syntax-and-semantics
% Section: valuations-sat

\documentclass[../../../include/open-logic-section]{subfiles}

\begin{document}

\olfileid{pl}{syn}{val}

\olsection{\usetoken{P}{valuation} અને સંતોષ}

\begin{defn}[!!^{valuation}s]
$\{\True,\False\}$ને ``સત્ય'' અને ``અસત્ય'' એ બે સત્યમૂલ્યોનો ગણ
માનો. $\Lang{L_0}$ માટેની \emph{!!{valuation}} એ ભાષાનાં
!!{propositional variable}sને $\True$ અથવા $\False$ નિયુક્ત કરતું
વિધેય~$\pAssign{v}$ છે; એટલે કે, $\pAssign{v}\colon
\PVar\to\{\True,\False\}$.
\end{defn}

\begin{defn}
  !!a{valuation}~$\pAssign{v}$ આપેલી હોય, ત્યારે મૂલ્યાંકન વિધેય
  $\pValue{v}\colon\Frm[L_0]\to\{\True,\False\}$ને નીચે પ્રમાણે
  અનુમાનાત્મક રીતે વ્યાખ્યાયિત કરો:
  \begin{align*}
    \iftag{prvFalse}{\pValue{v}(\lfalse) & = \False; \\}{}
    \iftag{prvTrue}{\pValue{v}(\ltrue) & = \True; \\}{}
    \pValue{v}(\Obj p_n) & = \pAssign{v}(\Obj p_n); \\
    \iftag{prvNot}{\pValue{v}(\lnot !A) & = \begin{cases}
        \True & \text{જો } \pValue{v}(!A) = \False;\\
        \False & \text{અન્યથા.}
      \end{cases} \\ }{}
    \iftag{prvAnd}{\pValue{v}(!A \land !B) & = \begin{cases}
        \True &
        \text{જો $\pValue{v}(!A) = \True$ અને $\pValue{v}(!B) = \True$;}\\
        \False &
        \text{જો $\pValue{v}(!A) = \False$ અથવા $\pValue{v}(!B) = \False$}.
      \end{cases}\\}{}
    \iftag{prvOr}{\pValue{v}(!A \lor !B) & = \begin{cases}
        \True &
        \text{જો $\pValue{v}(!A) = \True$ અથવા $\pValue{v}(!B) = \True$;}\\
        \False &
        \text{જો $\pValue{v}(!A) = \False$ અને $\pValue{v}(!B) = \False$}.
      \end{cases}\\}{}
    \iftag{prvIf}{\pValue{v}(!A \lif !B) & = \begin{cases}
        \True &
        \text{જો $\pValue{v}(!A) = \False$ અથવા $\pValue{v}(!B) = \True$;}\\
        \False &
        \text{જો $\pValue{v}(!A) = \True$ અને $\pValue{v}(!B) = \False$}.
      \end{cases}\\}{}
    \iftag{prvIff}{\pValue{v}(!A \liff !B) & = \begin{cases}
        \True &
        \text{જો $\pValue{v}(!A) = \pValue{v}(!B)$;}\\
        \False &
        \text{જો $\pValue{v}(!A) \neq \pValue{v}(!B)$}.
      \end{cases}}{}
\end{align*}
\end{defn}

\begin{explain}
આ ખંડો નીચેની સત્યાર્થતા સારણીઓને અનુરૂપ છે:
\begin{center}
\iftag{prvNot}{
  \begin{tabular}{|c||c|} \hline
    $!A$ & $ \lnot !A$ \\
    \hline \hline
    $\True$ & $\False$ \\
    $\False$ & $\True$ \\
    \hline
  \end{tabular}
}{}
\iftag{prvAnd}{
  \begin{tabular}{|cc||c|} \hline
    $!A$ & $!B$ & $!A \land !B$ \\
    \hline \hline
    $\True$ & $\True$ & $\True$ \\
    $\True$ & $\False$ & $\False$ \\
    $\False$ & $\True$ & $\False$ \\
    $\False$ & $\False$ & $\False$ \\
    \hline
  \end{tabular}
}{}
\iftag{prvOr}{
  \begin{tabular}{|cc||c|} \hline
    $!A$ & $!B$ & $!A \lor !B$ \\
    \hline \hline
    $\True$ & $\True$ & $\True$ \\
    $\True$ & $\False$ & $\True$ \\
    $\False$ & $\True$ & $\True$ \\
    $\False$ & $\False$ & $\False$ \\
    \hline
  \end{tabular}
}{}%
\iftag{notprvNot,notprvAnd,notprvOr,notprvIf}{}{\\[1em]}
\iftag{prvIf}{
  \begin{tabular}{|cc||c|} \hline
    $!A$ & $!B$ & $!A \lif !B$ \\
    \hline \hline
    $\True$ & $\True$ & $\True$ \\
    $\True$ & $\False$ & $\False$ \\
    $\False$ & $\True$ & $\True$ \\
    $\False$ & $\False$ & $\True$ \\
    \hline
  \end{tabular}
}{}
\iftag{prvIff}{
  \begin{tabular}{|cc||c|} \hline
    $!A$ & $!B$ & $!A \liff !B$ \\
    \hline \hline
    $\True$ & $\True$ & $\True$ \\
    $\True$ & $\False$ & $\False$ \\
    $\False$ & $\True$ & $\False$ \\
    $\False$ & $\False$ & $\True$ \\
    \hline
  \end{tabular}
}{}
\end{center}
\end{explain}

\begin{prob}
$\Lang{L_0}$માં ત્રિસ્થાની સંયોજક $\diamondsuit$ ઉમેરવાનું વિચારો,
જેનું મૂલ્યાંકન આ પ્રમાણે આપેલું છે:
\begin{gather*}
  \pValue{v}(\diamondsuit( !A , !B , !C ) ) = \begin{cases}
    \pValue{v}( !B ) &
    \text{જો $\pValue{v}(!A) = \True$;}\\
    \pValue{v}( !C ) &
    \text{જો $\pValue{v}(!A) = \False $}.
  \end{cases}
\end{gather*}
આ સંયોજકની સત્યાર્થતા સારણી લખો.
\end{prob}

\begin{thm}[સ્થાનિક નિર્ધારણ]
\ollabel{thm:LocalDetermination} ધારો કે $\pAssign{v_1}$ અને
$\pAssign{v_2}$ એ !!{valuation}s છે અને $!A$માં આવતાં
!!{propositional
variable}s પર સંમત થાય છે; એટલે કે, કોઈ
!!{formula}~$!A$માં $\Obj p_n$ આવે ત્યારે
$\pAssign{v_1}(\Obj p_n)=\pAssign{v_2}(\Obj p_n)$. ત્યારે
$\pValue{v_1}$ અને $\pValue{v_2}$ પણ~$!A$ પર સંમત થાય છે; એટલે કે,
$\pValue{v_1}(!A)=\pValue{v_2}(!A)$.
\end{thm}

\begin{proof}
$!A$ પર અનુમાનથી.
\end{proof}

\begin{defn}[સંતોષ]
\ollabel{defn:satisfaction} !!a{valuation}~$\pAssign{v}$ વડે
  !!a{formula}~$!A$ના \emph{સંતોષ}, $\pSat{v}{!A}$,નો ખ્યાલ આપણે
  નીચે પ્રમાણે અનુમાનાત્મક રીતે વ્યાખ્યાયિત કરી શકીએ છીએ.
  (``$\pSat{v}{!A}$ નથી'' એ કહેવા આપણે $\pSat/{v}{!A}$ લખીએ છીએ.)
\begin{enumerate}
\tagitem{prvFalse}{%
  \indcase{!A}{\lfalse}{$\pSat/{v}{\indfrm}$.}}{}

\tagitem{prvTrue}{%
  \indcase{!A}{\ltrue}{$\pSat{v}{\indfrm}$.}}{}

\item \indcase{!A}{\Obj p_i}{$\pSat{v}{\indfrm}$ ત્યારે જ જ્યારે
  $\pAssign{v}(\Obj p_i) = \True$.}

\tagitem{prvNot}{%
  \indcase{!A}{\lnot !B}{$\pSat{v}{\indfrm}$ ત્યારે જ જ્યારે
    $\pSat/{v}{!B}$.}}{}

\tagitem{prvAnd}{%
  \indcase{!A}{(!B \land !C)}{$\pSat{v}{\indfrm}$ ત્યારે જ જ્યારે
    $\pSat{v}{!B}$ અને $\pSat{v}{!C}$.}}{}

\tagitem{prvOr}{%
  \indcase{!A}{(!B \lor !C)}{$\pSat{v}{\indfrm}$ ત્યારે જ જ્યારે
    $\pSat{v}{!B}$ અથવા $\pSat{v}{!C}$ (અથવા બંને).}}{}

\tagitem{prvIf}{%
  \indcase{!A}{(!B \lif !C)}{$\pSat{v}{\indfrm}$ ત્યારે જ જ્યારે
    $\pSat/{v}{!B}$ અથવા $\pSat{v}{!C}$ (અથવા બંને).}}{}

\tagitem{prvIff}{%
  \indcase{!A}{(!B \liff !C)}{$\pSat{v}{\indfrm}$ ત્યારે જ જ્યારે
    $\pSat{v}{!B}$ અને $\pSat{v}{!C}$ બંને, અથવા $\pSat{v}{!B}$ અને
    $\pSat{v}{!C}$ એ બેમાંથી એકેય નહિ.}}{}
\end{enumerate}
જો $\Gamma$ એ !!{formula}sનો ગણ હોય, તો દરેક~$!A\in\Gamma$ માટે
$\pSat{v}{!A}$ હોય ત્યારે જ $\pSat{v}{\Gamma}$.
\end{defn}

\begin{prop}\ollabel{prop:sat-value}
  $\pSat{v}{!A}$ ત્યારે જ જ્યારે $\pValue{v}(!A)=\True$.
\end{prop}

\begin{proof}
  $!A$ પર અનુમાનથી.
\end{proof}

\begin{prob}
  \olref[pl][syn][val]{prop:sat-value} સાબિત કરો.
\end{prob}

\end{document}
